\documentclass[11pt]{article}

\usepackage[margin=1in]{geometry}
\usepackage{amsmath,amssymb}
\usepackage{booktabs}
\usepackage{hyperref}
\usepackage{enumitem}

\newcommand{\Fp}{\mathbb{F}_p}
\newcommand{\Z}{\mathbb{Z}}
\newcommand{\dx}{d_x}
\newcommand{\dy}{d_y}
\newcommand{\ellv}{\ell}

\title{Method Used by the Submitted secp256k1 Point-Addition Circuit}
\author{}
\date{}

\begin{document}
\maketitle

\section{Scope}

This document describes the circuit path used by the submission in
\texttt{src/point\_add/} when evaluated by \texttt{src/main.rs} through

\[
\texttt{cargo run --release}
\]

with no environment variables set.  Environment-gated diagnostics and
alternative experimental paths are intentionally ignored here.

The circuit implements generic affine addition on secp256k1,

\[
P=(x_P,y_P),\qquad Q=(x_Q,y_Q),\qquad R=P+Q,
\]

where \(P\) is quantum and \(Q\) is classical.  The evaluator avoids point
doubling, inverse-pair, and infinity cases, so the implemented path assumes the
affine denominators are nonzero.

\section{Evaluator Contract}

The evaluator in \texttt{src/main.rs} calls

\[
\texttt{point\_add::build()}
\]

and then analyzes the emitted operation stream.  The required interface is four
256-bit registers, declared in this order:

\[
\begin{array}{ccl}
0 &:& \texttt{target\_x} \text{, quantum qubits},\\
1 &:& \texttt{target\_y} \text{, quantum qubits},\\
2 &:& \texttt{offset\_x} \text{, classical bits},\\
3 &:& \texttt{offset\_y} \text{, classical bits}.
\end{array}
\]

The input is

\[
(\texttt{target\_x},\texttt{target\_y})=(x_P,y_P),
\qquad
(\texttt{offset\_x},\texttt{offset\_y})=(x_Q,y_Q),
\]

and the output must be

\[
(\texttt{target\_x},\texttt{target\_y})=(x_R,y_R).
\]

The operation stream is hashed with SHAKE256 and the hash-derived stream is used
both to generate test points and to seed the simulator.  Thus changing the
circuit changes the test inputs.  The evaluator generates random multiples of
the secp256k1 generator, filters out unsupported exceptional cases, runs the
circuit in 64-shot batches, and checks:

\begin{enumerate}[nosep]
  \item classical correctness against \texttt{WeierstrassEllipticCurve::add};
  \item zero global phase after the forward circuit;
  \item no non-register qubit garbage after the output registers are zeroed.
\end{enumerate}

The simulator counts \texttt{CCX} and \texttt{CCZ} as Toffoli gates.  It counts
\texttt{CX}, \texttt{CZ}, \texttt{Swap}, \texttt{Hmr}, and \texttt{R} as
Clifford gates.  \texttt{X} and \texttt{Z} are free in the reported metrics.

\section{Field and Curve Constants}

The circuit is specialized to secp256k1:

\[
p = 2^{256} - 2^{32} - 977,
\qquad
E: y^2=x^3+7 \pmod p.
\]

It uses the Solinas form

\[
2^{256} \equiv c \pmod p,
\qquad
c = 2^{32}+977
  = 2^{32}+2^{10}-2^6+2^4+1.
\]

This sparse representation is used throughout modular reduction.

\section{High-Level Point-Addition Formula}

The evaluator's affine addition formula is

\[
\lambda = \frac{y_Q-y_P}{x_Q-x_P},
\]

\[
x_R = \lambda^2 - x_P - x_Q,
\]

\[
y_R = \lambda(x_P-x_R)-y_P.
\]

The circuit first subtracts the classical offset from the quantum target:

\[
\dx = x_P-x_Q,
\qquad
\dy = y_P-y_Q.
\]

Since both numerator and denominator are negated,

\[
\lambda = \frac{\dy}{\dx}.
\]

The implementation stores the negative slope

\[
\ellv = -\lambda.
\]

\subsection{First inverse and slope computation}

The first Kaliski inversion is run on \(\dx\).  For the default no-env path, the
iteration count is

\[
t_1 = 404.
\]

The raw Kaliski output used by the circuit has the form

\[
I_1 = -\dx^{-1}2^{t_1} \pmod p.
\]

The circuit multiplies

\[
\ellv_{\mathrm{raw}} = \dy I_1
\]

and then applies \(t_1\) modular halvings:

\[
\ellv = 2^{-t_1}\ellv_{\mathrm{raw}}
      = -\frac{\dy}{\dx}
      = -\lambda.
\]

\subsection{Combined affine \(x\)- and \(y\)-update}

Instead of separately zeroing \(\dy\), computing \(x_R\), and then computing
\(y_R\), the default path uses a combined affine rearrangement.

Let

\[
A = \ellv^2 = \lambda^2.
\]

The circuit forms a temporary value

\[
B = A - 3x_Q.
\]

Then it updates the \(y\)-register as

\[
\dy \leftarrow \dy + \ellv B.
\]

Using \(\dy=\lambda(x_P-x_Q)\) and \(\ellv=-\lambda\),

\[
\dy + \ellv(A-3x_Q)
=
\lambda(x_P-x_Q)-\lambda(A-3x_Q)
=
\lambda(x_P+2x_Q-A).
\]

Since

\[
x_R = A-x_P-x_Q,
\]

we have

\[
y_R+y_Q
=
\lambda(x_P-x_R)-y_P+y_Q
=
\lambda(x_P+2x_Q-A).
\]

Thus after the combined update,

\[
\texttt{target\_y} = y_R+y_Q.
\]

The \(x\)-register is updated as

\[
\dx \leftarrow -(\dx-A+3x_Q)
              = A-\dx-3x_Q
              = A-(x_P-x_Q)-3x_Q
              = x_R-x_Q.
\]

So after the combined affine block,

\[
\texttt{target\_x}=x_R-x_Q,
\qquad
\texttt{target\_y}=y_R+y_Q,
\qquad
\ellv=-\lambda.
\]

\subsection{Second inverse and slope cleanup}

The slope register must be uncomputed.  The second Kaliski inversion is run on

\[
D = x_R-x_Q.
\]

The default iteration count for this second inverse is

\[
t_2 = 400.
\]

The raw inverse has the form

\[
I_2 = -D^{-1}2^{t_2}.
\]

Before using it, the circuit doubles the slope register \(t_2\) times:

\[
\ellv \leftarrow 2^{t_2}\ellv.
\]

At this point

\[
\texttt{target\_y}=y_R+y_Q.
\]

From the affine formula,

\[
y_R+y_Q=\lambda(x_Q-x_R)=-\lambda(x_R-x_Q)=\ellv D.
\]

Therefore

\[
I_2(y_R+y_Q)
=
(-D^{-1}2^{t_2})(\ellv D)
=
-2^{t_2}\ellv.
\]

The circuit adds this product into the slope register, giving zero:

\[
2^{t_2}\ellv + I_2(y_R+y_Q)=0.
\]

It then subtracts the classical \(y_Q\) from the \(y\)-register and adds the
classical \(x_Q\) to the \(x\)-register:

\[
y_R+y_Q-y_Q=y_R,
\qquad
x_R-x_Q+x_Q=x_R.
\]

The slope register is then clean and can be freed.

\section{Kaliski Inversion Subroutine}

The default path uses a qrisp-style binary Kaliski almost-inverse circuit.  It
does not overwrite the input denominator.  Instead it allocates a persistent
state

\[
(u,v,r,s,m_{\mathrm{hist}},f),
\]

where \(u,v,r,s\) are 256-bit quantum registers, \(m_{\mathrm{hist}}\) is an
iteration-history register, and \(f\) is a one-qubit active/terminal flag.

Initialization is

\[
u=p,
\qquad
v=x,
\qquad
r=0,
\qquad
s=1,
\qquad
f=1.
\]

Each iteration computes local branch flags, conditionally swaps \((u,v)\) and
\((r,s)\), conditionally performs

\[
v \leftarrow v-u,
\qquad
s \leftarrow s+r,
\]

then halves \(v\) and doubles \(r\).  The swaps are then undone, and the local
flags are uncomputed.  The branch history is stored in \(m_{\mathrm{hist}}\), so
the forward computation can later be explicitly reversed.

After the chosen fixed number of iterations, for the denominators exercised by
the submitted circuit, the state exposes

\[
r = -x^{-1}2^t \pmod p,
\]

with \(t=404\) in the first inverse and \(t=400\) in the second inverse.

The default implementation includes several important optimizations:

\begin{enumerate}
  \item \textbf{Bulk prefix iterations.}
  Early Kaliski iterations are known, for the tested nonzero denominators, to be
  nonterminal.  The default path uses specialized bulk iterations for this
  prefix.  The no-env caps are 378 bulk iterations for the first inverse and
  377 for the second inverse.

  \item \textbf{Truncated late-iteration widths.}
  The implementation uses bit-length bounds to shrink comparisons, swaps, and
  add/subtract slices in late iterations.  Roughly, after iteration \(i\), high
  bits above the bound \(2n-i\) are known zero in the denominator state.

  \item \textbf{Small-\(r\) modular doubling shortcut.}
  For iterations below the threshold 262, the top bit of \(r\) is known zero,
  so modular doubling of \(r\) needs no Solinas correction and is emitted as a
  plain shift.

  \item \textbf{Measurement-based uncomputation.}
  ANDs, carry chains, OR chains, and temporary products are often uncomputed by
  \texttt{Hmr} followed by classically-conditioned phase correction.  This
  removes many reverse Toffolis while preserving phase cleanliness.

  \item \textbf{State lifetime reduction.}
  After forward Kaliski and before the multiplication body, registers known to
  be zero or known constants are freed and then reallocated/reloaded for the
  backward pass.  In the default path this reduces the live qubit peak while
  preserving a clean explicit backward uncomputation.
\end{enumerate}

\section{Modular Arithmetic}

All arithmetic is over \(\Fp\), with \(n=256\).

\subsection{Cuccaro adders with measurement-based uncompute}

The primitive adder is a Cuccaro ripple-carry adder.  The default fast variant
computes carries using Toffolis and uncomputes many carry bits using
measurement-based reset:

\[
\texttt{Hmr}(t,m),
\qquad
\texttt{CZ}(a,b)\ \text{conditioned on }m.
\]

This is the standard measurement-based AND-uncompute pattern.  The evaluator's
global phase check is what enforces that all such measurement feedback is
phase-correct.

\subsection{Modular addition}

For \(a,b\in[0,p)\), the circuit computes \(a+b\bmod p\) using an extended
257-bit register.  Let

\[
c=2^{256}-p=2^{32}+977.
\]

First the circuit computes the unreduced sum \(S=a+b\).  It then adds \(c\).
The top bit of \(S+c\) indicates whether \(S\ge p\).  If reduction was not
needed, the added \(c\) is subtracted back.  If reduction was needed, the top
bit is cleared, yielding \(S-p\).  The reduction flag is then uncomputed by a
comparison identity.

\subsection{Modular subtraction}

For \(a-b\bmod p\), the circuit performs an extended subtraction and records the
borrow.  If there was an underflow, the wrapped low 256 bits are corrected using
the Solinas identity \(p\equiv -c \pmod{2^{256}}\).  The borrow flag is then
uncomputed through a comparison after temporarily negating the subtrahend.

\subsection{Doubling and halving}

Modular doubling is implemented as a left shift.  The old top bit becomes an
overflow flag.  If the overflow flag is set, the sparse constant \(c\) is added.

Modular halving is emitted as the inverse operation: the low bit determines the
correction, then a right shift is performed.  These operations are used heavily
to absorb the \(2^t\) scaling factors produced by Kaliski inversion.

\section{Multiplication and Reduction}

\subsection{Wide product generation}

The default path uses reversible wide-product generation followed by Solinas
reduction and then Bennett uncomputation of the wide product.

The main schoolbook product primitive is a Litinski-style controlled
add-subtract multiplier.  For \(n=256\), it uses a \(2n+1\)-bit accumulator and
performs \(n\) controlled add-subtract rows.  Corrections are then applied so
that the final shifted accumulator contains \(xy\).  The wide product is later
uncomputed by the exact inverse routine.

\subsection{Karatsuba where enabled by default}

With no environment variables, Karatsuba is enabled for the main pair-1
write-into-zero multiplication and for the pair-2 cleanup multiplication.

The 1-level Karatsuba split is

\[
x=x_0+2^{128}x_1,
\qquad
 y=y_0+2^{128}y_1,
\]

with

\[
z_0=x_0y_0,
\qquad
z_2=x_1y_1,
\]

\[
z_1=(x_0+x_1)(y_0+y_1)-z_0-z_2.
\]

The three half-size products are merged into the 512-bit temporary product.
After Solinas reduction into the target accumulator, the Karatsuba product is
run backward to clean the temporary registers.

\subsection{Symmetric squaring}

The combined affine block needs \(\ellv^2\).  The default path uses a symmetric
schoolbook squarer instead of a full multiplier.  Cross terms \(x_i x_j\) for
\(i<j\) are generated once and placed with the appropriate factor of two,
reducing the Toffoli count of the squaring step.

\subsection{Solinas reduction of a 512-bit product}

For a wide product

\[
T=T_0+2^{256}T_1,
\]

the Solinas prime gives

\[
T \equiv T_0 + cT_1 \pmod p.
\]

Using

\[
c=2^{32}+2^{10}-2^6+2^4+1,
\]

the circuit adds

\[
T_0
+T_1
+2^4T_1
-2^6T_1
+2^{10}T_1
+2^{32}T_1
\]

modulo \(p\).  The high half \(T_1\) is walked through the required shifts by
modular doubling and by a specialized 22-bit shift from position 10 to position
32.  After the reduction additions are complete, \(T_1\) is shifted/halved back
to its original value so that the wide product can be uncomputed.

\section{Default No-Environment Circuit Schedule}

The no-env path is:

\begin{enumerate}
  \item Allocate and declare the four interface registers.

  \item Compute
  \[
  \dx=x_P-x_Q,
  \qquad
  \dy=y_P-y_Q.
  \]

  \item Run the first raw Kaliski inverse on \(\dx\) with \(t_1=404\).

  \item Compute
  \[
  \ellv = \dy(-\dx^{-1}2^{404})2^{-404}
        = -\lambda.
  \]

  The multiplication here uses the default Karatsuba write-into-zero path.

  \item Use the combined affine block:
  \[
  A=\ellv^2,
  \qquad
  B=A-3x_Q,
  \]
  \[
  \texttt{target\_y}\leftarrow \dy+\ellv B = y_R+y_Q,
  \]
  \[
  \texttt{target\_x}\leftarrow A-\dx-3x_Q = x_R-x_Q.
  \]

  \item Run the second raw Kaliski inverse on \(x_R-x_Q\) with \(t_2=400\).

  \item Double the slope register 400 times and add
  \[
  I_2(y_R+y_Q)
  \]
  into it, cleaning the slope register to zero.

  \item Subtract the classical \(y_Q\):
  \[
  y_R+y_Q\mapsto y_R.
  \]

  \item Add the classical \(x_Q\):
  \[
  x_R-x_Q\mapsto x_R.
  \]

  \item Free all non-output quantum scratch.
\end{enumerate}

\section{Measured No-Environment Result}

A no-env release run of the submitted path produced:

\begin{center}
\begin{tabular}{lr}
\toprule
Metric & Value \\
\midrule
Emitted operations & 28,677,221 \\
Qubits & 2,715 \\
Classical bits & 2,339,473 \\
Average executed Toffoli & 3,942,753.000 \\
Average executed Clifford & 22,627,103.228 \\
Total Toffoli over 9,024 shots & 35,579,403,072 \\
Total Clifford over 9,024 shots & 204,186,979,531 \\
\bottomrule
\end{tabular}
\end{center}

The same run reported:

\[
\text{classical mismatches}=0,
\qquad
\text{phase-garbage batches}=0,
\qquad
\text{ancilla-garbage batches}=0.
\]

The submission also runs an internal alternate-seed diagnostic on the default
path before the main evaluator simulation.  In the measured no-env run, all five
alternate seeds over 4,096 shots each had zero classical mismatches, zero phase
garbage, and zero ancilla garbage.

\section{Summary}

The result is obtained by an affine secp256k1 point-addition circuit specialized
to the Solinas prime \(p=2^{256}-2^{32}-977\).  Its main ingredients are:

\begin{itemize}
  \item two raw Kaliski almost-inversions, used in scaled negative form;
  \item a combined affine \(x/y\) rearrangement that computes \(x_R-x_Q\) and
        \(y_R+y_Q\) directly;
  \item reversible Litinski-style schoolbook and 1-level Karatsuba
        multiplications;
  \item sparse Solinas modular reduction using
        \(2^{256}\equiv 2^{32}+977\);
  \item measurement-based uncomputation of carries, ANDs, and OR chains;
  \item explicit Kaliski backward passes using stored branch history;
  \item aggressive scratch lifetime reduction between Kaliski forward, body,
        and backward phases.
\end{itemize}

The evaluator accepts the circuit because the output registers match classical
secp256k1 addition on the Fiat-Shamir-derived tests, and because all measurement
phases and all quantum scratch are clean at the end of the forward circuit.

\end{document}
